\section{精确编译、原始投影与支持查询}\label{sec:compilation}
候选接口界限制关系状态，却没有直接给出查询答案。编译器还必须把同一历史中的关系约束一起保存，并在事件退休后保留它所决定的原始文字。本节说明这一连接。关系自动机和编译后查询是已有技术基础~\cite{kater2023,darwiche2002}；本文的结构定理为它们提供覆盖所有候选、由协议参数控制的输入计划。

\begin{algorithm}[ht]
\caption{由已认证计划构造精确表示}\label{alg:compile}
\begin{algorithmic}[1]
\Require 固定的 $I,\Gamma_I,G_I,\pi_I$；全部原始变量归属表
\Ensure 接受路径 DAG $\DAG_I$，或资源预算未完成状态
\State 建立初始联合摘要；将永久原始常量作为公共起始标签
\For{计划中的每个组 $j$}
  \State 按计划引入端点，安装尚未安装的标记
  \For{本层每个联合状态 $z$ 及组 $j$ 的每个选项 $b$}
    \State 在 $z$ 的同一历史摘要上加入 $b$ 的关系和来源计数
    \If{出现禁止环、计数溢出或退休读不满足恰一来源}
      \State 丢弃该转移
    \Else
      \State 投影到下一层端口；按全部摘要和计数合并下一层状态
      \State 保留标有 $b$ 的原始文字的有向边
    \EndIf
  \EndFor
\EndFor
\State 删除不通向接受终点的节点，保留共同起点与终点
\end{algorithmic}
\end{algorithm}

\subsection{二元组与稀疏修改序}
对每个保留的候选 $w\in D(r)$，一个二元组选择 $\RF wr$ 或 $\neg\RF wr$。正选择加入 $w\mathrel\rf r$ 及反向边，负选择只留下标签。读 $r$ 有一个已选来源计数，取值为 $0$ 或 $1$；再次选中立即拒绝，在其最后候选处理完时要求计数为 $1$。已删候选的负文字进入固定标签。

对每地址的 $L$，保留 $\Cov(L)$，并对每个 $L$ 不可比较写对用一个二元组选择一个方向。令这些方向与固定覆盖边组成 $g$。若 $g$ 无环，则 $g^+$ 在该地址上是严格全序：原来可比较的写对由 $\Cov(L)^+=L$ 连接，原来不可比较的写对已直接定向，故任意不同写可比较；无环性保证反自反和反对称。反之，任何扩展 $L$ 的严格全序选择相应方向，得到无环 $g$，其闭包恰为这个全序。因此只需监视 $g$ 无环，同时用 $g^+$ 解释 $\co$。不同地址之间不加入 CO 边。已由 $L$ 决定的 CO 位仍以原始常量文字归属到路径。

\subsection{联合端口摘要与安全遗忘}\label{sec:ports}
取字母 $s$ 表示真实线程邻接及稀疏初始化边，$r$ 表示已选 RF，$\bar r$ 表示其反向，$g$ 表示 CO 生成边。RA 禁止游走的语言可写为
\begin{equation}\label{eq:language}
 (s\mid r)^+\,(r\mid g\mid\bar r g^+)^*.
\end{equation}
它表示一段非空 $\hb$，后接零段或多段 RF、CO 或 FR；闭合游走恰违反 $\irr(\hb;\eco^?)$。一个四状态部分 DFA 足以识别该语言。把事件图与 DFA 作乘积，并在每个事件的接受状态到起始状态加入特殊\emph{标记边}，则违反 RA 等价于存在恰好经过一条标记边的闭合游走。另用一个单状态监视器检查 $g$ 无环。附录~\ref{app:automaton} 给出完整转移表及等价证明。

对某个已经选择的处理前缀，在乘积图中令 $P$ 为仅使用普通边的反身传递闭包，$B_\star$ 为已经安装的标记边集合，令 $T=P;B_\star;P$。$T$ 表示\emph{恰好一条}标记边，而不允许串接两个 $T$ 段。未来端口预先由全部尚未处理候选及尚未安装标记确定：$I_t$ 包含它们的目标端点，$O_t$ 包含它们的源端点。对于旧图，只需保存
\begin{equation}\label{eq:ports}
 P\cap(I_t\times O_t),\qquad T\cap(I_t\times O_t).
\end{equation}
彼此排斥的候选共同决定哪些端点将来可能被使用，但不会被同时加入已选前缀的关系。

\begin{lemma}[相同联合摘要具有相同后缀]\label{lem:ports}
在同一固定计划的同一层，若两个未被拒绝的前缀具有相同的全部监视器端口摘要及 RF 计数，则任一相同的未来二元选择序列，要么对二者都接受，要么对二者都拒绝。
\end{lemma}
\begin{proof}
固定未来选择。若完成后出现此前不存在的禁止闭合游走，它必使用新加入的边。沿该游走把旧边的极大连续段取出；每段的入口是某条未来边的目标，出口是某条未来边的源，故都位于式~\eqref{eq:ports} 保留的矩形中。对于 RA，整条游走恰有一个标记：若标记在旧段中，该段由 $T$ 表示，其余旧段由 $P$ 表示；若标记是新安装的边，所有旧段均由 $P$ 表示。因此用另一前缀的相同摘要替换各旧段，仍可连接同一批新边构成禁止游走。每段替代路径均存在于这一个替代前缀中，它们可串接成游走。CO 无环监视器使用同样论证，只需普通可达关系。

反方向交换两个前缀即可。来源计数相同且未来选项相同，溢出和最后恰一检查结果也相同。全部摘要和计数作为一个联合状态保存，因而结论同时作用于所有约束。若把不同前缀的不同监视器摘要独立合并，再任意组合，所需“同一个替代前缀”就会丢失，证明不再成立。\qedhere
\end{proof}

算法~\ref{alg:compile} 按固定计划构建分层 DAG。插入普通边或标记边时更新相应关系并立即检查禁止环；投影到下一层端口后，用整个联合状态作哈希键合并节点。附录~\ref{app:update} 给出具体增量公式。每条边记录本次选项的原始文字，合并节点不会合并或抹去不同入边上的文字。读源已处理完的事件可以退休，后续查询仍沿路径标签定位它的身份。



一个事件袋至多有 $K$ 个事件，RA 乘积图至多有 $4K$ 个顶点。它的普通与一标记矩阵最多用 $2(4K)^2$ 位；CO 矩阵最多用 $K^2$ 位；活跃读计数至多 $K$ 位。所以每层联合状态的粗上界为
\begin{equation}\label{eq:statebound}
 S_{\rm bd}=2^{33K^2+K}.
\end{equation}
实际端口矩形通常小于这些方阵，且大部分位模式不可达。实验中的 $S_{\max}$ 是实际达到的最大层状态数，含义不同于 $S_{\rm bd}$。

\subsection{对全部原始赋值的精确投影}\label{sec:projection}
对路径 DAG 可以用辅助变量表示一条从起点到终点的可用路径，得到 CNF $A_I(X,Y)$。附录~\ref{app:supportcnf} 给出线性大小且能传播原始支持的具体构造。以下结论中的 $x$ 遍历 $X$ 的\emph{所有}布尔赋值，包括违反来源恰一或 CO 传递性的赋值。

\begin{theorem}[精确原始投影]\label{thm:projection}
设 $I$ 为第~\ref{sec:problem}节的固定观测，$\Gamma_I$ 的每条删域及关系事实均可靠，所有值集合与分支限制已在候选域／常量中精确落实。设 $\pi_I$ 满足事件生存期、全部二元组覆盖及原始变量唯一归属条件。算法~\ref{alg:compile} 完成后，对所有 $x\in\{0,1\}^X$ 都有
\begin{equation}\label{eq:projection}
 \Phi_I(x)\quad\Longleftrightarrow\quad\exists y\ A_I(x,y).
\end{equation}
因此同一个 $\DAG_I$ 与 $A_I$ 对任意后续假设 $\alpha$ 精确保留全部原始 RF／CO 支持。
\end{theorem}
\begin{proof}
需要同时满足四项条件。

第一，观测和过滤精确。可靠性保证每个合法完成使用的来源仍在 $D$ 内，并满足全部固定文字；输入的值与守卫则确实通过这些域限制执行。对缓冲，先证明自身索引读唯一，再把原两读守卫化为对端来源的取值限制，这使化简在这个程序上是等价的。一般多读守卫若不能如此精确落实，就不在该前端接受范围内。仅有 $D$ 的上近似可靠性，不能代替这项反向保证。

第二，选择表示精确。二元 RF 组和计数检查保证每读恰一来源；逐地址的稀疏 CO 构造与全部扩展 $L$ 的严格全序一一对应。因此任何合法 $x$ 都给出计划中一条完整选择序列，任何通过这些局部检查的序列都定义满足原始来源和写序格式的赋值。

第三，关系检查和状态合并精确。完整 RF／CO 选择经式~\eqref{eq:language} 的监视器接受，当且仅当它满足 RA。引理~\ref{lem:ports} 保证退休和合并保持每个前缀的全部可接受后缀。沿接受图的一条路径逐层归纳：当前实际选择前缀属于该联合状态，图中下一条转移的可行性可由同后缀等价转移到这个前缀。因此不能通过接合不兼容历史制造新完成。

第四，原始身份没有在过滤或合并时丢失。每个原始位恰由一个固定文字或一个选择组赋值；初始化相关位、被删 RF 的负位、固定 CO 和已退休事件的位均包含在内。接受路径的完整标签正好给出一个 $x$。于是合法 $x$ 的序列不会被拒绝，产生接受路径；反向每条接受路径的标签满足输入观测和 RA。路径 CNF 的存在辅助赋值恰表示这种接受路径，得到式~\eqref{eq:projection}。\qedhere
\end{proof}

共同证书不能永久吸收可撤销的 $\alpha$。例如第~\ref{sec:example}节的晚发布假设使某个 CO 必然成立，但撤销它后两个方向应恢复。定理先构造只依赖 $I$ 的表示，再允许对 $X$ 任意条件化，从而保证这类恢复。

对于大小为 $K$ 的二元计划，规范化候选图的边数和组数为 $O(Kn)$：在每个事件的引入位置，它至多与袋中 $K-1$ 个旧端点关联，各候选对贡献常数个组。故图、标签及支持 CNF 的大小为
\begin{equation}\label{eq:size}
 O\bigl(|X|+Kn S_{\rm bd}\bigr),
\end{equation}
构造时间为 $T_{\rm pre}+O(|X|+KnS_{\rm bd}\operatorname{poly}(K))$ 的单元操作。二进制编码事件和变量标识时另有对数因子。预处理及其位费用详见附录~\ref{app:complexity}。即使 $K$ 固定，$|X|$ 仍可能为 $\Theta(n^2)$，所以显式回答所有原始位的成本仍然增长。

\subsection{完整支持与按需日志诊断}\label{sec:queries}
给定文字合取 $\alpha$，临时禁用标签与 $\alpha$ 矛盾的边。在其余图中分别计算从起点可达和能够到达终点的节点。若终点不可达，返回 $\UNSAT$；否则，一个原始文字 $\lambda$ 有支持，当且仅当某条携带 $\lambda$ 的边同时位于这两类节点之间。每个这样的边都在一条完整接受路径上，反向每条支持见证路径也会经过对应文字。扫描费用线性于图、标签、假设和输出长度。撤销文字后在永久图上重新扫描即可。支持 CNF 通过单位传播实现相同的原始文字排除；该结论只针对原始变量条件化，不要求对辅助变量任意赋值也传播完备。

完整支持可以服务一个具体的来源日志任务。客户端依次加入 RF／CO 字段，发现矛盾时按固定顺序尝试删除字段，只要余下记录仍不可满足就保留删除，最终得到删除极小冲突。若要枚举全部单字段替换，对每个字段 $f$ 暂时移除其记录，在其余记录下计算该字段的支持，输出与旧值不同的全部支持值。替换可行恰等价于它在其余记录下有支持，所以这种枚举完整；它不要求对所有字段联合替换进行搜索。

在缓冲实例中，记录 $P1\_peer$ 来自 $A_1$，同时记录 $Q0\_peer$ 来自 $W_2$，两项分别可行而合取冲突。完整单字段修复包括把前者改为初始化 $A_0$，或把后者改为 $W_1$。客户端可以接受一项修复，再撤销并恢复记录，按要求输出完整支持快照。这个任务利用了精确投影与撤销能力；删除检查和单字段修复本身使用常规算法。第~\ref{sec:diagnosis}节允许 SAT 基线仅回答当时需要的一致性或字段查询，以衡量完整编译在实际请求结构中的收益。
